% Generated draft: requires mathematical and source review.
\documentclass[11pt]{article}
\usepackage[margin=25mm]{geometry}
\usepackage{fontspec}
\setmainfont{Latin Modern Roman}
\usepackage{amsmath,amssymb,mathtools,unicode-math}
\setmathfont{Latin Modern Math}
\usepackage{enumitem,needspace,xurl,hyperref}
\setlength{\parindent}{0pt}
\setlength{\parskip}{5pt}
\setlength{\emergencystretch}{3em}
\newcommand{\N}{\mathbb{N}}
\newcommand{\Z}{\mathbb{Z}}
\newcommand{\Q}{\mathbb{Q}}
\newcommand{\R}{\mathbb{R}}
\newcommand{\C}{\mathbb{C}}
\newcommand{\F}{\mathbb{F}}
\newcommand{\A}{\mathbb{A}}
\newcommand{\PP}{\mathbb{P}}
\newcommand{\E}{\mathbb{E}}
\newcommand{\eps}{\varepsilon}
\newcommand{\dd}{\,\mathrm{d}}
\begin{document}
\section*{UnsolvedMath ID: 701058}
\textbf{Source problem code:} catalogue entry
\subsection{Definitions and mathematical statement}
Let \(\Gamma\) be a finite graph, write \(V(\Gamma)\) for its vertex set, and put \(G=\operatorname{Aut}(\Gamma)\). For \(\alpha\in V(\Gamma)\), let
\[
\Gamma(\alpha)=\{\beta\in V(\Gamma):\beta\text{ is adjacent to }\alpha\},\qquad
G_\alpha=\{g\in G:g(\alpha)=\alpha\}.
\]
The graph is vertex-transitive if \(G\) acts transitively on \(V(\Gamma)\); its valency is then the common neighbourhood size \(v=|\Gamma(\alpha)|\). The supplied record does not specify graph conventions or whether connectedness is implicit. No connectedness hypothesis has been added here; this omission requires clarification before the formulation can be considered definitive.

For a permutation group \(L\leq\operatorname{Sym}(\Omega)\) on a nonempty set \(\Omega\), call \(L\) quasiprimitive if it is transitive and every nonidentity normal subgroup of \(L\) is transitive on \(\Omega\). Define the induced local permutation group by
\[
L_\alpha=\{g|_{\Gamma(\alpha)}:g\in G_\alpha\}
\leq\operatorname{Sym}(\Gamma(\alpha)).
\]
Using the standard induced-action interpretation of the source terminology, \(\Gamma\) is locally quasiprimitive if \(L_\alpha\) is quasiprimitive for every vertex \(\alpha\). Vertex transitivity makes it sufficient to check one vertex. This condition concerns the induced permutation group, whereas the conjectured bound concerns the entire stabilizer \(G_\alpha\), including the kernel of its action on \(\Gamma(\alpha)\).

\textbf{Conjecture (C. E. Praeger, 1999; source formulation).} There exists a function \(f:\N\to\N\) such that, for every finite vertex-transitive, locally quasiprimitive graph \(\Gamma\) of valency \(v\), and every \(\alpha\in V(\Gamma)\),
\[
|\operatorname{Aut}(\Gamma)_\alpha|\leq f(v).
\]
Here \(\N\) retains the source's unspecified convention. The source does not explain whether valency zero is included or how quasiprimitivity on an empty neighbourhood should be interpreted. These boundary conventions, together with the missing connectedness convention, remain unresolved.

\textbf{Supplied reductions and qualifications.} The record states that proving the conjecture can be reduced to graphs whose full automorphism group satisfies
\[
\text{every }1\ne N\trianglelefteq G\text{ has at most two orbits on }V(\Gamma).
\]
It attributes this reduction to C. E. Praeger, \emph{Ars Combin.} \textbf{19 A} (1985), 149--163.

The source also describes the analogous conjecture for locally primitive graphs, attributed to R. Weiss in 1978, as still open. Here a transitive permutation group is primitive if it preserves no partition other than the singleton partition and the one-part partition. For non-bipartite graphs, the source reports a ``reduction'' of Weiss' conjecture to the case of an almost simple full automorphism group, citing M. Conder, C. H. Li and C. E. Praeger (2000). A finite group \(H\) is almost simple if \(T\leq H\leq\operatorname{Aut}(T)\) for some nonabelian finite simple group \(T\), identifying \(T\) with its inner automorphism group. The source refers to Problem 14.46 for definitions, but that record is not supplied.

The supplied \emph{Comment of 2013} reports that these conjectures have been proved when there is an upper bound on the degree of any alternating group occurring as a quotient of a subgroup, citing Praeger, Pyber, Spiga and Szab\'{o} (2012). An alternating group of degree \(m\) is \(A_m\); occurrence as a quotient of a subgroup means a section \(K/J\cong A_m\), with \(J\trianglelefteq K\). The supplied wording does not explicitly identify the ambient group containing \(K\), the uniform scope of the degree bound, or the permitted dependence of the resulting stabilizer bound. Accordingly, this comment is retained as a qualified source report rather than expanded into a precise additional theorem.

The supplied \emph{Comment of 2021} reports that the conjecture has been proved when \(L_\alpha\) has an abelian regular normal subgroup, citing P. Spiga (2016). This means that there is an abelian subgroup \(A\trianglelefteq L_\alpha\) whose action on \(\Gamma(\alpha)\) is transitive and whose point stabilizers are trivial. All reductions, special-case results and open-status descriptions above are supplied source assessments, not independently verified literature conclusions.
\subsection{Short English statement}
Praeger's conjecture asks whether the size of a vertex stabilizer in the full automorphism group of a finite vertex-transitive, locally quasiprimitive graph is bounded solely by the valency. The supplied record labels the problem open and reports reductions and special cases. Connectedness, graph conventions and the valency-zero boundary require clarification; the reported literature assessments have not been independently verified.
\subsection{Sources}
\begin{itemize}
\item Supplied local UnsolvedMath registry record, permanent ID \(701058\), exact source code \texttt{catalogue entry}, title \emph{Kourovka Notebook Problem 14.73}. Proposed by C. E. Praeger in 1999. The registry labels the problem \texttt{open}; this is supplied metadata, not an independently verified mathematical assessment. Upstream revision: \texttt{d371ee917ac9647fbb06cd6fc529deec6f3dda82}; supplied record update: \texttt{2026-10-16T14:00:00Z}.
\item \emph{The Kourovka Notebook (Unsolved Problems in Group Theory)}, Problem 14.73, originating in the 14th issue (1999), as represented in the supplied 21st-issue edition (2026), arXiv:1401.0300v48. The supplied collection description identifies E. I. Khukhro and V. D. Mazurov as editors. The statement cross-references Problem 14.46, whose contents are not supplied.
\item C. E. Praeger, \emph{Ars Combin.} \textbf{19 A} (1985), 149--163. Supplied citation for the reduction to full automorphism groups whose nonidentity normal subgroups have at most two vertex orbits; no article title or URL is supplied.
\item R. Weiss, 1978: supplied attribution for the analogous locally primitive conjecture. No full bibliographic citation is supplied.
\item M. Conder, C. H. Li and C. E. Praeger, \emph{Proc. Edinburgh Math. Soc. (2)} \textbf{43}, no. 1 (2000), 129--138. Supplied citation for the reported almost-simple reduction of Weiss' conjecture for non-bipartite graphs; no article title or URL is supplied.
\item C. E. Praeger, L. Pyber, P. Spiga and E. Szab\'{o}, \emph{Proc. Amer. Math. Soc.} \textbf{140}, no. 7 (2012), 2307--2318. Cited in the supplied \emph{Comment of 2013} concerning bounded alternating-section degrees. This is supplied dated commentary, not independently verified here.
\item P. Spiga, \emph{Bull. London Math. Soc.} \textbf{48}, no. 1 (2016), 12--18. Cited in the supplied \emph{Comment of 2021} concerning an abelian regular normal subgroup in the induced local group. This is supplied dated commentary, not independently verified here.
\end{itemize}
\end{document}
